% Fragmento propuesto para insertar despues de la Figura~\ref{fig:app_factor_effects_paired}.

\subsection*{Lectura de los efectos emparejados}

La Figura~\ref{fig:app_factor_effects_paired} compara el efecto medio de distintas condiciones experimentales manteniendo emparejadas el resto de configuraciones. Los deltas son, en general, pequeños respecto a su variabilidad, lo que indica que ningún factor aislado domina de forma estable el resultado. En mejora del \emph{incumbent}, la ausencia de ruido muestra un efecto positivo moderado frente a ruido gaussiano, mientras que Sobol frente a aleatorio y ARD frente a kernel base presentan efectos débiles.

En MAE, ARD tiende a mejorar ligeramente el resultado medio, pero con alta dispersión. El muestreo Sobol no muestra una ventaja clara frente al aleatorio y el efecto del ruido sobre el MAE es pequeño. Por tanto, la figura sugiere que el comportamiento final depende más del benchmark y de la interacción entre condiciones que de una única decisión experimental aislada.
